% ----------------------------------------------------------
\chapter{Introdução}\label{cap:introducao}
% ----------------------------------------------------------

A constante matemática \texorpdfstring{$\pi$}{pi} é definida como a razão entre o comprimento de qualquer circunferência e seu respectivo diâmetro em um espaço euclidiano, constituindo um dos números fundamentais da matemática \cite{diniz2018}. Seu valor aproximado é 3,14159, sendo classificado como um número irracional e transcendente, pois não pode ser expresso como razão entre inteiros nem como raiz de qualquer polinômio com coeficientes inteiros.

O estudo de \texorpdfstring{$\pi$}{pi} remonta à Antiguidade, quando civilizações como a egípcia e a babilônica já empregavam aproximações numéricas em problemas de natureza geométrica. Posteriormente, Arquimedes de Siracusa desenvolveu métodos baseados em polígonos inscritos e circunscritos para estabelecer limites superiores e inferiores para seu valor, inaugurando abordagens rigorosas de estimativa. Na matemática moderna, \texorpdfstring{$\pi$}{pi} surge naturalmente em diversas áreas, incluindo trigonometria, análise matemática, séries infinitas e integrais definidas.

\section{Processos Estocásticos}

Processos estocásticos \cite{gardiner2009} constituem uma classe de modelos matemáticos empregados na descrição de fenômenos cuja evolução é influenciada por componentes aleatórios. Formalmente, um processo estocástico é definido como uma família de variáveis aleatórias indexadas por um conjunto, que geralmente representa o tempo ou uma variável espacial.

Do ponto de vista aplicado, processos estocásticos são amplamente utilizados na modelagem de sistemas dinâmicos sujeitos a incertezas e flutuações aleatórias. Assim, constituem ferramentas essenciais na análise quantitativa de sistemas complexos, possibilitando previsões probabilísticas e a compreensão do comportamento médio.

\section{Objetivos}

\subsection{Objetivo Geral}
Estimar a constante matemática \texorpdfstring{$\pi$}{pi} utilizando o método numérico estocástico de Monte Carlo e realizar uma análise rigorosa do erro algorítmico associado.

\subsection{Objetivos Específicos}
\begin{alineas}
    \item Apresentar a fundamentação teórica sobre o Método de Monte Carlo e a geometria da circunferência unitária;
    \item Implementar um algoritmo em linguagem Python para a simulação de dardos aleatórios;
    \item Calcular o erro absoluto, relativo e relativo percentual a cada passo da iteração;
    \item Avaliar a convergência do método à medida que o tamanho da amostra aumenta.
\end{alineas}